\chapter{Tests}

\section{Tests fonctionnels}
Le backend est couvert par une suite de tests automatisés (\texttt{pytest},
\texttt{pytest-asyncio}) exécutée contre une base SQLite en mémoire, couvrant notamment~:
la création et le cycle de vie des projets, l'amorçage automatique du Knowledge Graph,
le déroulement complet d'un entretien, la génération des exigences, l'analyse de
sécurité, le moteur de revue et la génération documentaire multi-formats.

\TODO{Indiquer le nombre exact de tests et le taux de couverture au moment de la
rédaction finale du rapport (à régénérer via \texttt{pytest --cov} juste avant la
remise).}

\section{Tests d'intégration}
Les tests d'API exercent les endpoints HTTP de bout en bout via un client de test
asynchrone, garantissant que les routes, les schémas Pydantic et les services métier
restent cohérents entre eux.

\section{Tests de sécurité}
Le moteur de sécurité est testé sur des cas explicites~: absence de fabrication d'une
base légale de traitement des données non confirmée, cohérence entre les acteurs
déclarés dans les exigences et la matrice RBAC générée, et non-duplication d'une
analyse de sécurité déjà exécutée pour un même projet.

\section{Tests de génération documentaire}
Le générateur de documents est testé pour garantir qu'un projet fraîchement créé (sans
réponse d'entretien) ne produit \textbf{aucune donnée fabriquée}~: dictionnaire de
données vide, absence d'analyse de sécurité tant qu'elle n'a pas été exécutée, et
mention explicite des sections volontairement non générées (diagrammes BPMN, maquettes
visuelles).

\section{Tests de l'assistant IA}
\TODO{Décrire les tests spécifiques à l'assistant IA (longueur des questions, respect
du périmètre du projet, gestion des erreurs de connexion au fournisseur) une fois
formalisés et automatisés.}
